% GENERATED FILE. Edit canonical lessons, then run npm run generate:books.
% canonical-source-hash: fnv1a64:0edc9f0f99ce6d1e
% canonical-lessons: HI-C455-gam, HI-C455-niras, HI-C455-mayus, HI-C455-kurhna, HI-C455-dukhana

\chapter{Sadness and Regret}
\label{ch:hi-sadness-and-regret}
\hlchaptermodality{455}

\begin{chapteropening}
\textbf{By the end of this chapter:} \emph{I can say that someone was let down or sorry.}

Complete the last lesson of chapter 455: I can say that someone was let down or sorry.
\end{chapteropening}

\section[\texorpdfstring{\hi{ग़म} (ġam)}{ग़म (ġam)}]{\hi{ग़म} (ġam) — grief; sorrow at a loss}

\label{lesson:HI-C455-gam}



\begin{quote}
Before the new one: say the Hindi for strange, then the Hindi for a marvel.
\end{quote}

\subsection*{You'll want to know: \hi{ग़म}}
\textbf{\hi{ग़म}} — \emph{ġam} — \textquotedblleft{}grief; sorrow at a loss\textquotedblright{}.

\begin{quote}
\hi{पूरा} \hi{परिवार} \hi{दादा} \hi{की} \hi{मौत} \hi{के} \hi{ग़म} \hi{में} \hi{है।}
\end{quote}

\begin{quote}
\emph{pūrā parivār dādā kī maut ke ġam meṁ hai.}
\end{quote}

\begin{quote}
The whole family is grieving over my grandfather's death.
\end{quote}

Masculine.

\begin{cousinweb}
\textbf{\hi{ग़म}} (\emph{ġam}) is usually traced to Arabic \emph{ġamm}, \textquotedblleft{}grief, distress\textquotedblright{}, which reached Hindi through Persian.
\end{cousinweb}

\begin{grammarlens}[title={What you've built}]
One more word for describing what you lived and felt.
\end{grammarlens}

\subsection*{Your turn}
\begin{itemize}
\raggedright
  \item \emph{Say it:} \emph{ġam}
  \item \emph{Say it:} \emph{ġam}, once more
  \item \emph{Say it:} \emph{pūrā parivār dādā kī maut ke ġam meṁ hai.} — The whole family is grieving over my grandfather's death.
  \item \emph{Recall:} say the Hindi for strange, then the Hindi for a marvel, then say \emph{ġam} again
\end{itemize}

\begin{tcolorbox}[breakable,colback=teal!4,colframe=teal!35!black,title={Before you move on}]
What does \hi{ग़म} mean? (\textquotedblleft{}Grief; sorrow at a loss\textquotedblright{}.) Say it once more.
\end{tcolorbox}
\section[\texorpdfstring{\hi{निराश} (nirāś)}{निराश (nirāś)}]{\hi{निराश} (nirāś) — disappointed; without hope}

\label{lesson:HI-C455-niras}



\begin{quote}
Before the new one: say the Hindi for a marvel, then the Hindi for grief.
\end{quote}

\subsection*{You'll want to know: \hi{निराश}}
\textbf{\hi{निराश}} — \emph{nirāś} — \textquotedblleft{}disappointed; without hope\textquotedblright{}.

\begin{quote}
\hi{परीक्षा} \hi{का} \hi{नतीजा} \hi{देखकर} \hi{अरुण} \hi{निराश} \hi{हुआ।}
\end{quote}

\begin{quote}
\emph{parīkṣā kā natījā dekhkar Aruṇ nirāś huā.}
\end{quote}

\begin{quote}
Arun was disappointed when he saw his exam result.
\end{quote}

\begin{cousinweb}
\textbf{\hi{निराश}} (\emph{nirāś}) is a tatsama, borrowed from Sanskrit \emph{nirāśa}, \textquotedblleft{}without any hope or wish\textquotedblright{} (Monier-Williams).
\end{cousinweb}

\begin{grammarlens}[title={What you've built}]
One more word for describing what you lived and felt.
\end{grammarlens}

\subsection*{Your turn}
\begin{itemize}
\raggedright
  \item \emph{Say it:} \emph{nirāś}
  \item \emph{Say it:} \emph{nirāś}, once more
  \item \emph{Say it:} \emph{parīkṣā kā natījā dekhkar Aruṇ nirāś huā.} — Arun was disappointed when he saw his exam result.
  \item \emph{Recall:} say the Hindi for a marvel, then the Hindi for grief, then say \emph{nirāś} again
\end{itemize}

\begin{tcolorbox}[breakable,colback=teal!4,colframe=teal!35!black,title={Before you move on}]
What does \hi{निराश} mean? (\textquotedblleft{}Disappointed; without hope\textquotedblright{}.) Say it once more.
\end{tcolorbox}
\section[\texorpdfstring{\hi{मायूस} (māyūs)}{मायूस (māyūs)}]{\hi{मायूस} (māyūs) — dejected; downcast}

\label{lesson:HI-C455-mayus}



\begin{quote}
Before the new one: say the Hindi for grief, then the Hindi for disappointed.
\end{quote}

\subsection*{You'll want to know: \hi{मायूस}}
\textbf{\hi{मायूस}} — \emph{māyūs} — \textquotedblleft{}dejected; downcast\textquotedblright{}.

\begin{quote}
\hi{टीम} \hi{हार} \hi{गई}, \hi{इसलिए} \hi{सब} \hi{खिलाड़ी} \hi{मायूस} \hi{थे।}
\end{quote}

\begin{quote}
\emph{ṭīm hār gaī, isliye sab khilāṛī māyūs the.}
\end{quote}

\begin{quote}
The team lost, so all the players were downcast.
\end{quote}

\begin{cousinweb}
\textbf{\hi{मायूस}} (\emph{māyūs}) is usually traced to Arabic \emph{maʾyūs}, \textquotedblleft{}one given up as hopeless\textquotedblright{}, which reached Hindi through Persian.
\end{cousinweb}

\begin{grammarlens}[title={What you've built}]
One more word for describing what you lived and felt.
\end{grammarlens}

\subsection*{Your turn}
\begin{itemize}
\raggedright
  \item \emph{Say it:} \emph{māyūs}
  \item \emph{Say it:} \emph{māyūs}, once more
  \item \emph{Say it:} \emph{ṭīm hār gaī, isliye sab khilāṛī māyūs the.} — The team lost, so all the players were downcast.
  \item \emph{Recall:} say the Hindi for grief, then the Hindi for disappointed, then say \emph{māyūs} again
\end{itemize}

\begin{tcolorbox}[breakable,colback=teal!4,colframe=teal!35!black,title={Before you move on}]
What does \hi{मायूस} mean? (\textquotedblleft{}Dejected; downcast\textquotedblright{}.) Say it once more.
\end{tcolorbox}
\section[\texorpdfstring{\hi{कुढ़ना} (kuṛhnā)}{कुढ़ना (kuṛhnā)}]{\hi{कुढ़ना} (kuṛhnā) — to fret; to grieve inwardly}

\label{lesson:HI-C455-kurhna}



\begin{quote}
Before the new one: say the Hindi for disappointed, then the Hindi for dejected.
\end{quote}

\subsection*{You'll want to know: \hi{कुढ़ना}}
\textbf{\hi{कुढ़ना}} — \emph{kuṛhnā} — \textquotedblleft{}to fret; to grieve inwardly\textquotedblright{}.

\begin{quote}
\hi{अरुण} \hi{अपनी} \hi{हार} \hi{पर} \hi{बहुत} \hi{कुढ़ा।}
\end{quote}

\begin{quote}
\emph{Aruṇ apnī hār par bahut kuṛhā.}
\end{quote}

\begin{quote}
Arun fretted bitterly over his defeat.
\end{quote}

Past (he …): \hi{वह} \hi{बहुत} \hi{कुढ़ा}, \emph{vah bahut kuṛhā}.

\begin{cousinweb}
\textbf{\hi{कुढ़ना}} (\emph{kuṛhnā}) is a tadbhava: CDIAL lists Hindi \emph{kuṛhnā}, \textquotedblleft{}to be annoyed, to grieve\textquotedblright{}, under Sanskrit \emph{kruddhá}, \textquotedblleft{}angry\textquotedblright{} (CDIAL 3598). It is the anger or sorrow one swallows instead of showing.
\end{cousinweb}

\begin{grammarlens}[title={What you've built}]
One more word for describing what you lived and felt.
\end{grammarlens}

\subsection*{Your turn}
\begin{itemize}
\raggedright
  \item \emph{Say it:} \emph{kuṛhnā}
  \item \emph{Say it:} \emph{kuṛhnā}, once more
  \item \emph{Say it:} \emph{Aruṇ apnī hār par bahut kuṛhā.} — Arun fretted bitterly over his defeat.
  \item \emph{Recall:} say the Hindi for disappointed, then the Hindi for dejected, then say \emph{kuṛhnā} again
\end{itemize}

\begin{tcolorbox}[breakable,colback=teal!4,colframe=teal!35!black,title={Before you move on}]
What does \hi{कुढ़ना} mean? (\textquotedblleft{}To fret; to grieve inwardly\textquotedblright{}.) Say it once more.
\end{tcolorbox}
\section[\texorpdfstring{\hi{दुखाना} (dukhānā)}{दुखाना (dukhānā)}]{\hi{दुखाना} (dukhānā) — to hurt someone's feelings}

\label{lesson:HI-C455-dukhana}



\begin{quote}
Before the new one: say the Hindi for dejected, then the Hindi for to fret.
\end{quote}

\subsection*{You'll want to know: \hi{दुखाना}}
\textbf{\hi{दुखाना}} — \emph{dukhānā} — \textquotedblleft{}to hurt someone's feelings\textquotedblright{}.

\begin{quote}
\hi{मीरा} \hi{ने} \hi{कड़वी} \hi{बात} \hi{कहकर} \hi{अपनी} \hi{माँ} \hi{का} \hi{दिल} \hi{दुखाया।}
\end{quote}

\begin{quote}
\emph{Mīrā ne kaṛvī bāt kahkar apnī māṁ kā dil dukhāyā.}
\end{quote}

\begin{quote}
Mira hurt her mother's feelings with a harsh remark.
\end{quote}

Past, with \hi{ने} (he or she …): \hi{उसने} \hi{दिल} \hi{दुखाया}, \emph{usne dil dukhāyā}.

\begin{cousinweb}
\textbf{\hi{दुखाना}} (\emph{dukhānā}) is a tadbhava: CDIAL lists Hindi \emph{dukhānā}, the causative, beside \emph{dukhnā}, \textquotedblleft{}to ache\textquotedblright{}, under Sanskrit \emph{duḥkhati}, \textquotedblleft{}pains\textquotedblright{} (CDIAL 6376). \emph{dil dukhānā}, \textquotedblleft{}to make the heart ache\textquotedblright{}, is the usual phrase.
\end{cousinweb}

\begin{grammarlens}[title={What you've built}]
One more word for describing what you lived and felt.
\end{grammarlens}

\subsection*{Your turn}
\begin{itemize}
\raggedright
  \item \emph{Say it:} \emph{dukhānā}
  \item \emph{Say it:} \emph{dukhānā}, once more
  \item \emph{Say it:} \emph{Mīrā ne kaṛvī bāt kahkar apnī māṁ kā dil dukhāyā.} — Mira hurt her mother's feelings with a harsh remark.
  \item \emph{Recall:} say the Hindi for dejected, then the Hindi for to fret, then say \emph{dukhānā} again
\end{itemize}

\begin{tcolorbox}[breakable,colback=teal!4,colframe=teal!35!black,title={Before you move on}]
What does \hi{दुखाना} mean? (\textquotedblleft{}To hurt someone's feelings\textquotedblright{}.) Say it once more.
\end{tcolorbox}
